\sectionkicker{Source-backed technical notes}
\subsection*{Agent Safety \& Security — 記憶・retrieval・harnessが新しい攻撃面になる: Technical Notes}
本欄は記事本文で圧縮した一次資料上の情報を、比較・再検証しやすい形で整理したものである。時系列、確認済みの主張、留意点、一次資料URLを掲載する。
\medskip
\subsection*{Theme at a glance}
\addcontentsline{toc}{subsection}{Theme at a glance}
\begin{center}
\footnotesize
\begin{tabularx}{\linewidth}{@{}p{0.20\linewidth}p{0.10\linewidth}p{0.14\linewidth}X@{}}
\toprule
資料 & 位置づけ & 種別 & 時系列 \\
\midrule
From Prompt Injection to Persistent Control: Defending Agentic Harness Against Trojan Backdoors & 主要資料 & 論文 & 2026-05-29 (論文\_RELEASE) \\
Hijacking Agent Memory: Stealthy Trojan Attacks Through Conversational Interaction & 主要資料 & 論文 & 2026-05-28 (論文\_RELEASE) \\
Relevance as a Vulnerability: How Web Retrieval Degrades Safety Alignment in LLM Agents & 補足資料 & 論文 & 2026-05-28 (論文\_RELEASE) \\
Send a SCOUT First: Pre-hoc Reasoning for Adaptive Detector Allocation in Prompt-Injection Defense & 補足資料 & 論文 & 2026-05-29 (論文\_RELEASE) \\
\bottomrule
\end{tabularx}
\end{center}
\normalsize

\begin{technicalnote}{From Prompt Injection to Persistent Control: Defending Agentic Harness Against Trojan Backdoors}{主要資料}
\begin{tabularx}{\linewidth}{@{}>{\bfseries}p{0.22\linewidth}X@{}}
組織 & - \\
種別 & 論文 \\
時系列 & 2026-05-29 (論文\_RELEASE) \\
\end{tabularx}
\smallskip
{\bfseries 一次資料から整理したtechnical points}
\begin{itemize}[leftmargin=1.5em,itemsep=0.35em]
\item \textbf{一次情報で確認できる事実}: The preserved primary source identifies “From Prompt Injection to Persistent Control: Defending Agentic Harness Against Trojan Backdoors” on 2026-05-29.
\item \textbf{Author claim}: ClawTrojan models multi-step persistent prompt-injection attacks in local agentic harnesses, while DASGuard traces and removes untrusted control-like workspace content.
\end{itemize}
{\bfseries 読む際の境界}
\begin{itemize}[leftmargin=1.5em,itemsep=0.35em]
\item \textbf{分析上の留意点}: The preserved original arXiv metadata/abstract is primary-paper evidence, but experimental results and generalization are author-reported and were not independently reproduced in this Evidence review.
\end{itemize}
{\bfseries 一次資料}
\begingroup\sloppy
\begin{itemize}[leftmargin=1.5em,itemsep=0.25em]
\item {\scriptsize\url{http://arxiv.org/abs/2605.31042v1}}
\end{itemize}
\endgroup
\end{technicalnote}
\medskip

\begin{technicalnote}{Hijacking Agent Memory: Stealthy Trojan Attacks Through Conversational Interaction}{主要資料}
\begin{tabularx}{\linewidth}{@{}>{\bfseries}p{0.22\linewidth}X@{}}
組織 & - \\
種別 & 論文 \\
時系列 & 2026-05-28 (論文\_RELEASE) \\
\end{tabularx}
\smallskip
{\bfseries 一次資料から整理したtechnical points}
\begin{itemize}[leftmargin=1.5em,itemsep=0.35em]
\item \textbf{一次情報で確認できる事実}: The preserved primary source identifies “Hijacking Agent Memory: Stealthy Trojan Attacks Through Conversational Interaction” on 2026-05-28.
\item \textbf{Author claim}: MemPoison attacks selective long-term memory pipelines through conversational interactions using trigger/payload binding, entity masquerading, and embedding optimization.
\end{itemize}
{\bfseries 読む際の境界}
\begin{itemize}[leftmargin=1.5em,itemsep=0.35em]
\item \textbf{分析上の留意点}: The preserved original arXiv metadata/abstract is primary-paper evidence, but experimental results and generalization are author-reported and were not independently reproduced in this Evidence review.
\end{itemize}
{\bfseries 一次資料}
\begingroup\sloppy
\begin{itemize}[leftmargin=1.5em,itemsep=0.25em]
\item {\scriptsize\url{http://arxiv.org/abs/2605.29960v1}}
\end{itemize}
\endgroup
\end{technicalnote}
\medskip

\begin{technicalnote}{Relevance as a Vulnerability: How Web Retrieval Degrades Safety Alignment in LLM Agents}{補足資料}
\begin{tabularx}{\linewidth}{@{}>{\bfseries}p{0.22\linewidth}X@{}}
組織 & - \\
種別 & 論文 \\
時系列 & 2026-05-28 (論文\_RELEASE) \\
\end{tabularx}
\smallskip
{\bfseries 一次資料から整理したtechnical points}
\begin{itemize}[leftmargin=1.5em,itemsep=0.35em]
\item \textbf{一次情報で確認できる事実}: The preserved primary source identifies “Relevance as a Vulnerability: How Web Retrieval Degrades Safety Alignment in LLM Agents” on 2026-05-28.
\item \textbf{Author claim}: AgentREVEAL analyzes how web retrieval integration and retrieved-content relevance can degrade safety alignment in LLM agents.
\end{itemize}
{\bfseries 読む際の境界}
\begin{itemize}[leftmargin=1.5em,itemsep=0.35em]
\item \textbf{分析上の留意点}: The preserved original arXiv metadata/abstract is primary-paper evidence, but experimental results and generalization are author-reported and were not independently reproduced in this Evidence review.
\end{itemize}
{\bfseries 一次資料}
\begingroup\sloppy
\begin{itemize}[leftmargin=1.5em,itemsep=0.25em]
\item {\scriptsize\url{http://arxiv.org/abs/2605.29224v1}}
\end{itemize}
\endgroup
\end{technicalnote}
\medskip

\begin{technicalnote}{Send a SCOUT First: Pre-hoc Reasoning for Adaptive Detector Allocation in Prompt-Injection Defense}{補足資料}
\begin{tabularx}{\linewidth}{@{}>{\bfseries}p{0.22\linewidth}X@{}}
組織 & - \\
種別 & 論文 \\
時系列 & 2026-05-29 (論文\_RELEASE) \\
\end{tabularx}
\smallskip
{\bfseries 一次資料から整理したtechnical points}
\begin{itemize}[leftmargin=1.5em,itemsep=0.35em]
\item \textbf{一次情報で確認できる事実}: The preserved primary source identifies “Send a SCOUT First: Pre-hoc Reasoning for Adaptive Detector Allocation in Prompt-Injection Defense” on 2026-05-29.
\item \textbf{Author claim}: SCOUT dynamically allocates heterogeneous prompt-injection detectors and optional LLM judging based on predicted per-request reliability and latency.
\end{itemize}
{\bfseries 読む際の境界}
\begin{itemize}[leftmargin=1.5em,itemsep=0.35em]
\item \textbf{分析上の留意点}: The preserved original arXiv metadata/abstract is primary-paper evidence, but experimental results and generalization are author-reported and were not independently reproduced in this Evidence review.
\end{itemize}
{\bfseries 一次資料}
\begingroup\sloppy
\begin{itemize}[leftmargin=1.5em,itemsep=0.25em]
\item {\scriptsize\url{http://arxiv.org/abs/2605.30837v2}}
\end{itemize}
\endgroup
\end{technicalnote}
\medskip
